\documentclass{resume}
\usepackage{zh_CN-Adobefonts_external} 
\usepackage{linespacing_fix}
\usepackage{cite}
\begin{document}
\pagenumbering{gobble}

%**********************************姓名********************************************
\name{杨弘毅}
%**********************************联系信息****************************************
\contactInfo{(+86) 请填写手机号}{yhy520@email.swu.edu.cn}
%**********************************其他信息****************************************
\otherInfo{性别：男}{籍贯：请填写}{}{}
\otherInfo{GitHub：HongyiYoung.github.io}{个人主页：Vinyyang.github.io}{}{}
%*********************************照片**********************************************
% 如需照片请取消注释并添加照片到images文件夹，命名为you.jpg
% \yourphoto{0.15}

%**********************************教育背景**********************************************
\section{教育背景}

\datedsubsection{\textbf{西南大学}，计算机科学与技术专业，\textit{本科在读}}{2023.09 - 2027.06}
\begin{itemize} [parsep=1ex]
  \item \textbf{学业成绩}：GPA 3.0/5.0，专业排名 20/123（前10.63\%），获国家奖学金（2024）
  \item \textbf{英语水平}：CET-4 451分，CET-6 470分
  \item \textbf{核心课程}：数据结构与算法、计算机视觉、机器学习、嵌入式系统设计
  \item \textbf{研究方向}：计算机视觉、自然语言处理、具身AI系统
\end{itemize}

%**********************************科研项目**********************************************
\section{科研项目}

\datedsubsection{\textbf{自动驾驶感知系统研发}，算法工程师实习生}{2024.7 -- 2024.8}
在自动驾驶快速发展背景下，我负责开发2D语义分割和3D目标检测算法，构建从数据集标注到实车部署的完整感知系统。基于MMSegmentation框架实现像素级语义分割，支持19类道路场景要素识别；利用MMDetection3D处理LiDAR点云数据，实现多模态融合的3D目标检测；通过IsaacSim仿真平台进行算法验证和参数优化。最终实现像素级分割精度>85\%，3D检测距离>60米，系统在城市、高速、恶劣天气等多场景下稳定运行。

\textbf{实习单位}：重庆中科汽车软件创新中心

%**********************************竞赛获奖情况**********************************************
\section{竞赛获奖情况}

\datedsubsection{\textbf{全国大学生电子设计竞赛}，省级一等奖}{2025.8}
在72小时限时竞赛中，我担任技术负责人，主导设计了集成计算机视觉、嵌入式控制和精密机械的自主瞄准系统。基于OpenMV H7平台开发透视自适应目标检测算法，设计对角线交点法计算精确中心坐标，开发实时圆形轨迹生成算法实现车辆运动与绘图轨迹同步控制。最终实现<1.5cm瞄准精度，6cm±0.3cm圆形绘制精度，同步误差<1/4周期，获得省级一等奖。

\datedsubsection{\textbf{美国大学生数学建模竞赛（MCM）}，Finalist}{2025.1}
针对"过度旅游：走向衰落还是可持续发展？"问题，我担任队长负责数学模型构建与优化。建立了可持续旅游发展模型（STD模型）和模糊-CRITIC评估模型（FCE模型），运用微分方程、逻辑方程和直觉模糊集等方法，构建游客数量、环境质量和资本投资的动态平衡系统。通过敏感性分析和多城市适应性验证，为朱诺市旅游委员会提供了可持续发展建议。最终获得Finalist（决赛入围奖），该奖项比例<2\%。

\datedsubsection{\textbf{全国大学生数学建模竞赛}，省级二等奖}{2024.11}
针对"某乡村作物最优种植方案的决策研究"问题，我主要负责数学模型的算法实现与结果分析。基于动态规划思想构建混合整数规划模型，采用模拟退火算法和蒙特卡洛方法处理不确定性因素，通过层次聚类分析作物间的替代性和互补性。建立风险评估模型，为2024-2030年种植策略提供最优方案，实现7年总收益最大化。获得省级二等奖。

%**********************************技能特长**********************************************
\section{技能特长}

\begin{itemize}[parsep=0.5ex]
  \item \textbf{编程语言}：Python（熟练）、C/C++（熟练）、R（精通）、LaTeX（精通）
  \item \textbf{专业技能}：计算机视觉、机器学习、深度学习框架（PyTorch、MMSegmentation、MMDetection3D）
  \item \textbf{开发工具}：STM32开发（Keil5、HAL库）、嵌入式系统、ROS、NVIDIA IsaacSim
  \item \textbf{语言能力}：中文（母语）、粤语（母语）、英语（流利，CET-4 451，CET-6 470）
\end{itemize}

%**********************************自我评价**********************************************
\section{自我评价}

\begin{itemize}[parsep=0.5ex]
  \item 具备扎实的计算机科学理论基础和丰富的工程实践经验，在计算机视觉和嵌入式系统领域有深入研究
  \item 拥有优秀的问题分析和解决能力，能够将理论知识有效转化为实际应用，具备从概念设计到产品部署的全流程开发经验
  \item 具有良好的团队协作精神和项目管理能力，在多个竞赛项目中担任技术负责人，能够协调团队高效完成复杂任务
  \item 学习能力强，对新技术保持敏锐嗅觉，持续关注人工智能、自动驾驶、具身智能等前沿技术发展
\end{itemize}

\end{document}
